\documentclass[10pt,a4paper]{amsart}
\usepackage[margin=19mm]{geometry}
\usepackage{amsmath,amssymb,mathtools}
\usepackage[scr=rsfs,cal=euler]{mathalfa}
\usepackage{xfrac}
\usepackage[unicode,hidelinks,bookmarks=false]{hyperref}
\newcommand{\ie}{i.e.\ }
\newcommand{\cf}{cf.\ }
\newcommand{\resp}{respectively\ }
\newcommand{\Subsection}{\subsection}
\providecommand{\nobreakdash}{\nobreak\hbox{-}}
\setlength{\emergencystretch}{2em}
\multlinegap0pt
\usepackage{bm}
\usepackage{bbm}
\usepackage{dsfont}
\usepackage{xspace}
\usepackage{cancel}
\usepackage{mathtools}
\usepackage{amsthm}
\makeatletter
\@ifundefined{theorem}{\newtheorem{theorem}{Theorem}}{}
\@ifundefined{lemma}{\newtheorem{lemma}[theorem]{Lemma}}{}
\@ifundefined{prop}{\newtheorem{prop}[theorem]{Proposition}}{}
\@ifundefined{proposition}{\newtheorem{proposition}[theorem]{Proposition}}{}
\makeatother
\DeclareMathOperator{\Tr}{{\rm Tr}}
\newcommand{\E}{{\mathbf E}}
\newcommand{\RR}{{\mathbb R}}
\newcommand{\ZZ}{{\mathbb Z}}
\newcommand{\tr}{{\rm Tr}}
\title[Subordinacy theory for long-range operators: hyperbolic geod]{Subordinacy theory for long-range operators: hyperbolic geodesic flow insights and monotonicity theory}
\author{Zhenfu Wang, Disheng Xu, Qi Zhou}
\date{Source: arXiv:2506.23098v1 (CC BY 4.0)}
\begin{document}
\maketitle

\noindent\emph{Source and licence.} Subordinacy theory for long-range operators: hyperbolic geodesic flow insights and monotonicity theory. Version arXiv:2506.23098v1 (CC BY 4.0), distributed under the Creative Commons Attribution 4.0 International licence, as recorded in the arXiv licence metadata for this exact version and shown on the abstract page https://arxiv.org/abs/2506.23098v1. Full rights notice: licenses/arxiv-2506.23098-CC-BY-4.0.txt.\par\smallskip

\noindent\textit{Editorial scope.} Counted unit: the Lemma whose source label is \texttt{imm} in the article (source0.tex lines 1893--1900, source label \texttt{imm}), delivered together with the article's own complete original proof of it (source0.tex lines 1901--1965, one \texttt{proof} environment of 65 lines). No mathematics is written, repaired or re-derived: statement and proof are the article's own text line for line. Required context retained uncounted: lem: key ang est1 (lines 1261--1268); angle1 (lines 1263--1265); lem: ang imp slow1 (lines 1289--1300); prop: variation center (lines 1485--1491). Every definition, display and label of those lines is retained unchanged. Omitted from this package and not redistributed: 1--1260, 1266--1288, 1301--1484, 1492--1892, 1966--3448. Nothing inside a retained excerpt is omitted or altered: every retained body is delivered exactly as the article's lines are written, so no omission receipt of any kind accompanies this package. Citations: the retained excerpts carry no citation command at all, so no bibliography entry of the article is transcribed and no reference is redistributed. Reference adaptation: none was needed and none was made; every reference inside the delivered unit resolves inside it, so no unresolved reference or citation is produced. Printed slips kept literal: no printed word, symbol, hypothesis or number of the counted statement or of its proof was altered. Layout and class: the article's own class is \texttt{amsart} with packages including inputenc, amsmath, amsthm, amsfonts, amssymb, epsfig, enumerate, color, hyperref, xy, float, tikz, graphicx, mathrsfs; the wrapper is the lane's pinned \texttt{amsart} editorial wrapper at 10pt on A4, which restates the theorem-like environments the retained text needs and adds the article's own macro definitions that the retained text needs (\texttt{\textbackslash E}, \texttt{\textbackslash RR}, \texttt{\textbackslash Tr}, \texttt{\textbackslash ZZ}, \texttt{\textbackslash tr}), with extra wrapper packages (bm, bbm, dsfont, xspace, cancel, mathtools). The delivered numbering is the unit's own. Assets: no retained span contains a figure environment, a graphics-inclusion command, a PDF-page inclusion, a table or a TikZ picture; no image file is copied into this package, the package declares no asset dependency and no image file is redistributed with it. Licence: version arXiv:2506.23098v1 of this article is distributed under the Creative Commons Attribution 4.0 International licence, as recorded in the arXiv licence metadata for this exact version and shown on the abstract page https://arxiv.org/abs/2506.23098v1. A full rights notice is delivered at licenses/arxiv-2506.23098-CC-BY-4.0.txt. Characters: every retained excerpt is pure ASCII, so the delivered engine, XeLaTeX with its default Latin Modern text font, reports no missing character.
\begin{proposition}\label{lem: key ang est1}
Assume that \( (T,A_E) \) is partially hyperbolic, and the energy \( E \in \mathcal{G}^s_\omega \). There exists an orthonormal basis \( \{\mathbf{u}_j(\omega)\}_{j=1}^m \) of \( \mathbb{C}^m(0,\omega) \) such that for any \( \{\mathbf{v}_{j,E+i\epsilon}(\omega)\}_{j=1}^m \in \mathbb{C}^m(1,\omega) \) that continuously depends on \( E+i\epsilon \) when \( \epsilon \neq 0 \), there exists a constant \( \gamma = \gamma(E,\omega,\mathbf{v}) > 0 \) satisfying the angular lower bound:
\begin{equation}\label{angle1}
\angle\left((\mathbf{u}_j(\omega), \mathbf{v}_{j,E+i\epsilon}(\omega)), E^s_{A_{E+i\epsilon}}(\omega)\right) = \gamma>0,
\end{equation}
for any \( 1 \leq j \leq m \) and \( \epsilon > 0 \).
Furthermore, $\gamma$ is continuously dependent on $E.$
\end{proposition}

\begin{lemma}\label{lem: ang imp slow1}
Suppose that the uniform angle condition \eqref{angle1} holds, and assume furthermore
\begin{equation*}
\left(\mathbf{u}_j(\omega), \mathbf{v}_{j,E+i\epsilon}(\omega)\right) \in \mathcal{E}^s_{A_{E+i\epsilon}}(\omega).
\end{equation*}
There exists \( {C}_1 = {C}_1(E,\omega,\mathbf{v}) > 0 \) such that for all \( k \in \mathbb{Z} \):
\begin{equation*}
\|(A_{E+i\epsilon})_k(\omega)(\mathbf{u}_j(\omega), \mathbf{v}_{j,E+i\epsilon}(\omega))^{T}\|
\geq {C}_1\|((A_{E+i\epsilon})_k(\omega)|_{E^c_{A_{E+i\epsilon}}})^{-1}\|^{-1}\|(\mathbf{u}_j(\omega), \mathbf{v}_{j,E+i\epsilon}(\omega))^{T}\|.
\end{equation*}
Furthermore, $C_1$ is continuously dependent on $E.$
\end{lemma}

\begin{prop}\label{prop: variation center}Suppose that the  Schr\"odinger cocycle  $(T, A_E)$ is partially hyperbolic for some $E\in \RR$. Then there exists constant $C_2=C_2(E)>1$ 
such that for any $\epsilon\in \RR$ sufficiently close to $0$, any $\omega\in \Omega$, any $n\in \ZZ^+$ we have 
\begin{equation*}\label{eqn: est center variation}
\| (A_{E+i\epsilon})_{ n} (\omega)|_{E^c_{A_{E+i\epsilon}}}\|, \| (A_{E+i\epsilon})_{ n} (\omega)|_{E^c_{A_{E+i\epsilon}}})^{-1}\|  \leq  C_1C(n)\exp (C_1C(n)\epsilon n),
\end{equation*}
 where  $C(n):=\max_{0\leq s\leq n}\|(A_E)_{s}(\omega)|_{E^c_{A_E}}\|^2$. Furthermore, $C_2$ is continuously dependent  on $E.$
\end{prop}

\begin{lemma}\label{imm}
    Let $E\in \mathcal{G}_\omega^s$, then 
    \begin{equation}\label{eqn: criterion key ineq}
     \tr\Im M_{E+i\epsilon}^+(\omega)\geq 
	 C_5^{-1}(m+\tr C^{-1} M_{E+i\epsilon}^+(\omega)(M_{E+i\epsilon}^+(\omega))^\ast(C^{-1})^*),   
    \end{equation}
	 where $C_5:=C_4\max_{0\leq s\leq 3[\epsilon^{-1}]}\|(A_E)_{s}(\omega)|_{E^c_{A_E}}\|^6$, $C_4$ is a constant dependent only on $E,\omega$, and is continuously dependent on $E.$
\end{lemma}

\begin{proof}


For each basis vector $\mathbf{u}_j(\omega)$ given by Proposition \ref{lem: key ang est1}, consider the unique $\ell^2(\mathbb{Z}_+,\mathbb{C}^m)$ solution:
\begin{equation*}
	H_{V,T,\omega}\mathbf{u}_{E+i\epsilon}(n,\omega) = (E+i\epsilon)\cdot \mathbf{u}_{E+i\epsilon}(n,\omega)
\end{equation*}
satisfying the boundary condition  $\mathbf{u}_{E+i\epsilon}(0,\omega)=\mathbf{u}_j(\omega)$, such $\mathbf{u}_{E+i\epsilon}$ is unique and by the definition of $M^+_{E+i\epsilon}$, we have that $\mathbf{u}_{E+i\epsilon}(1,\omega)=-C^{-1}M_{E+i\epsilon}^+\mathbf{u}_j(\omega)$ which holomorphically depends on $E+i\epsilon$ when $\epsilon\neq 0$.

By the selection, 
$\left(\mathbf{u}_{E+i\epsilon}(0,\omega), \mathbf{u}_{E+i\epsilon}(1,\omega)\right)^{T} \in \E^s_{A_{E+i\epsilon}}(\omega),$
then by 
Lemma \ref{lem: ang imp slow1} and  Proposition \ref{prop: variation center}, we get that there exists $C_6=C_6(E,\omega,V)$ which is continuously dependent on $E$, such that 
\begin{eqnarray*}
&&\|(\mathbf{u}_{E+i\epsilon}(2k+1,\omega), \mathbf{u}_{E+i\epsilon}(2k+2,\omega))\|^2\\
&\geq& (C_5C(2k+1)\exp(C_5C(2k+1)\epsilon (2k+1))^{-2}\cdot \|(\mathbf{u}_{E+i\epsilon}(0,\omega), \mathbf{u}_{E+i\epsilon}(1,\omega))\|  \\
&\geq &(C_5C(2k+1)\exp(C_5C(2k+1)\epsilon (2k+1))^{-2} \cdot (1+\|C^{-1}M_{E+i\epsilon}^+(\omega)\cdot \mathbf{u}_j(\omega)\|^2).
\end{eqnarray*}
Hence using the increasing property of $C(n)$, we get for any $N$, 
\begin{eqnarray}
&&\langle\mathbf{u}_j(\omega), \Im M_{E+i\epsilon}^+(\omega)\cdot \mathbf{u}_j(\omega)\rangle \nonumber \\
&=&-\Im  \langle\mathbf{u}_{E+i\epsilon}(0,\omega),  C\mathbf{u}_{E+i\epsilon}(1,\omega)\rangle \nonumber \\
&=&\epsilon\cdot \sum_{k=0}^\infty \|(\mathbf{u}_{E+i\epsilon}(2k+1,\omega), \mathbf{u}_{E+i\epsilon}(2k+2,\omega))\|^2  \nonumber \\
&\geq &\epsilon\cdot \sum_{k=0}^N \|(\mathbf{u}_{E+i\epsilon}(2k+1,\omega), \mathbf{u}_{E+i\epsilon}(2k+2,\omega))\|^2 \nonumber\\
&\geq & \sum_{k=0}^N (C_5C(2N+1)\exp(C_5C(2N+1)\epsilon (2k+1))^{-2} \cdot (1+\|C^{-1}M_{E+i\epsilon}^+(\omega)\cdot \mathbf{u}_j(\omega)\|^2)\nonumber\\
&\geq & C_5^{-2}C(2N+1)^{-2}\cdot(\epsilon\exp(-2C_5C(2N+1)\epsilon)\cdot \frac{1-\exp(-4N C_5 C(2N+1)\epsilon)}{1-\exp(-4C_5C(2N+1)\epsilon)}\nonumber\\ && \cdot(1+\|C^{-1}M_{E+i\epsilon}^+(\omega)\cdot \mathbf{u}_j(\omega)\|^2). \label{eqn: est Im M M2 Simon}
\end{eqnarray}
 We know on $[0,1]$, $2xe^{-x}/(1-e^{-2x})$ has a lower bound $C>0$. Denote by $x(N):=2C_6C(2N+1)\epsilon$, then we have that the right hand side of \eqref{eqn: est Im M M2 Simon} is greater than 
\begin{eqnarray*}
\frac{1}{2}  C_6^{-3}C(2N+1)^{-3} x\frac{e^{-x}(1-e^{-2Nx})}{1-e^{-2x}} \cdot (1+\|C^{-1}M_{E+i\epsilon}^+(\omega)\cdot \mathbf{u}_j(\omega)\|^2) .
\end{eqnarray*}
Let $N_0=N_0(\epsilon)$ be the largest $N>0$ such that $x(N_0)\leq 1$, then when $N_0<\infty$, we have that $x(N_0)\geq x(N_0+1)\cdot \|A\|_{C^0}^{-4}\geq \|A\|_{C^0}^{-4}$. Moreover if we pick $\epsilon$ small enough (only depend on $\|A\|_{C^0}$), we can assume 
\begin{equation*}
N_0> \|A\|_{C^0}^4.
\end{equation*}
Then if $N_0<\infty$, we have 
\begin{equation}\label{eqn: est const}
 x(N_0)N_0\geq 1.
\end{equation}
If $N_0\leq \epsilon^{-1}$, take $N$ be $N_0$ in the right hand side of \eqref{eqn: est Im M M2 Simon}, by \eqref{eqn: est const} and $x(N)\leq 1$, we have that the right hand side of \eqref{eqn: est Im M M2 Simon} is greater than 
\begin{eqnarray}
&& \frac{1}{2}  C_6^{-3}C(2N_0+1)^{-3} \cdot (1+\|C^{-1}M_{E+i\epsilon}^+(\omega)\cdot \mathbf{u}_j(\omega)\|^2)\cdot x\frac{e^{-x}(1-e^{-2Nx})}{1-e^{-2x}} \nonumber\\
&\geq& \frac{1}{2}  C_6^{-3}C(2N_0+1)^{-3} \cdot (1+\|C^{-1}M_{E+i\epsilon}^+(\omega)\cdot \mathbf{u}_j(\omega)\|^2)\cdot 
\inf_{x\in[0,1]}\frac{xe^{-x}}{1-e^{-2x}}(1-e^{-2})\nonumber\\
&\geq &C_6 C(2N_0+1)^{-3}(1+\|C^{-1}M_{E+i\epsilon}^+(\omega)\cdot \mathbf{u}_j(\omega)\|^2)\ \nonumber\\
&\geq & C_6 C(3[\epsilon^{-1}])^{-3}(1+\|C^{-1}M_{E+i\epsilon}^+(\omega)\cdot \mathbf{u}_j(\omega)\|^2) .\label{eqn: last step simon est}
\end{eqnarray}
If $N_0>\epsilon^{-1}$, we take $N$ be $[\epsilon^{-1}]+1$, then $x(N)\leq 1$ and $x(N)N\geq 1$ (We can choose $C_6$ to be large, depending only on $E, \omega$, and continuously depending on $E$), we have that the right hand side of \eqref{eqn: est Im M M2 Simon} is greater than
$$\frac{1}{2}  C_6^{-3}C(2N+1)^{-3} \cdot (1+\|C^{-1}M_{E+i\epsilon}^+(\omega)\cdot \mathbf{u}_j(\omega)\|^2)\cdot 
\inf_{x\in[0,1]}\frac{xe^{-x}}{1-e^{-2x}}(1-e^{-2})\nonumber,$$
hence we have a similar estimate as \eqref{eqn: last step simon est}. In summary, we have 
\begin{equation}\label{eqn: est spec H+}
\langle\mathbf{u}_j(\omega), \Im M_{E+i\epsilon}^+(\omega)\cdot \mathbf{u}_j(\omega)\rangle \geq C_4 C(3[\epsilon^{-1}])^{-3}(1+\|C^{-1}M_{E+i\epsilon}^+(\omega)\cdot \mathbf{u}_j(\omega)\|^2).
\end{equation}


Let $\mathbf{u}_j$ run over the base we chose,  \eqref{eqn: est spec H+} implies the following inequality for positive definite matrices
\begin{equation*}
\begin{aligned}
	\tr\Im M_{E+i\epsilon}^+(\omega)&\geq  C_4C(3[\epsilon^{-1}])^3(m+\Tr(M_{E+i\epsilon}^+(\omega))^\ast(C^{-1})^*C^{-1} M_{E+i\epsilon}^+(\omega))\\
	&= C_4C(3[\epsilon^{-1}])^3(m+\tr C^{-1} M_{E+i\epsilon}^+(\omega)(M_{E+i\epsilon}^+(\omega))^\ast(C^{-1})^*),
\end{aligned}
\end{equation*}
where $C_5:=C_4C(3[\epsilon^{-1}])^3$.
\end{proof}

\end{document}
